\section{Introduction}\label{sec:intro}
Audio deepfake detectors are used to flag speech that may have been synthesized or voice converted. In an authentication workflow, a false alarm can divert a bona fide recording to manual review, whereas a missed spoof can leave impersonation undetected. The cost of either decision depends on the workflow and the prevalence of spoofed traffic. A detector chosen by aggregate discrimination alone can behave differently at a fixed operating point after a recording passes through ordinary signal processing.

Noise and lossy coding are particularly relevant because they can occur in either order. Speech recorded in a noisy room may be encoded before transmission, while background audio can also be introduced after decoding. These pipelines preserve the intended class of a source, but change the evidence presented to a detector. Even with identical noise and codec settings, reversing the operators can change the waveform, spectral features, ranking, and fixed-threshold decisions. Measurements across matched views can isolate the influence of the implemented signal-processing order within a recorded source.

The accepted LAVA conference manuscript benchmarks six heterogeneous detectors. The extension presented here addresses a different question through a matched intervention on one detector architecture. Within each of two initialization blocks, the training recordings, order, optimizer, numerical backend, and checkpoint rule are shared. The controlled factor is whether each training recording is presented clean or with a view sampled from a channel bank. Both fitted models are evaluated on each development recording and condition. This enables direct accounting of ranking changes, class-specific error changes, and which sources cross the same decision threshold.

The analysis asks whether aggregate robustness moves in the same direction in both blocks, which attacks and speakers bear the associated errors, and how those errors change under reversal of a conventional codec and noise chain. The resulting evidence is a controlled comparison and a source-linked diagnostic. Sections~\ref{sec:related}--\ref{sec:setup} specify the design, and Sections~\ref{sec:results}--\ref{sec:limits} distinguish measured findings from possible explanations and threats to validity.

\subsection{Operational question and scope}
The action following a detector decision determines which error is costly. A system that routes suspicious clips to an analyst incurs more review when authentic audio crosses the threshold. A system used to prioritize potentially manipulated evidence has a different cost for missing a synthetic or converted clip. A balanced validation cohort therefore answers a diagnostic question about conditional errors. It does not estimate workload for a deployment with an unknown spoof prevalence. The operational quantities reported here are denominated within class and condition, so that a practitioner can consider the two errors separately.

Clean evaluation provides a useful reference point because a detector can rank bona fide and spoofed recordings well before channel processing changes their score distributions. In audio production and communications, however, noise and compression can occur in either order. It is possible to keep the source, nominal noise level and codec settings fixed while varying the order. An observed score change then identifies sensitivity to the specified composition on that recording. Comparison across two different recording lists would entangle this effect with differences in speakers, content and synthesis examples.

The training comparison addresses a complementary question. Exposure to individual channel views may help on some transformed evaluation views even when the specific compositions were absent from the training-view bank. The experiment measures this transfer within ASVspoof 2019 LA under matched training and checkpoint selection. It records whether the rank benefit and fixed-threshold behavior agree, and whether the observed differences survive inspection at the attack and speaker level. This scope places claims about ordered signal processing and model response on separate evidential footing from claims about deployment or forensic cue preservation.

\subsection{Research questions}
The conference study compared available artifacts under a common offline test and timing contract. Its clean comparison used 2,737 checksum-disjoint recordings with paired correctness tests and stratified bootstrap intervals. Its 100-recording stress diagnostic summarized individual noise, codec, and simulated replay conditions. The six artifacts had heterogeneous training histories, including a warm-up EfficientNet-B0 checkpoint and externally pretrained RawNet2 and AASIST references. The journal uses a separate ASVspoof design to examine the effect of training-view policy and to link errors across two ordered transformations of each source. The historical ranking remains specific to the conference cohort.

The journal study asks three narrower questions. First, does changing only the training-view policy improve ranking under each registered stress condition for two matched initializations? Second, which class, speaker, and attack groups account for the difference between aggregate ranking and errors at the same operating point? Third, what can a same-source reversal of noise and conventional codec order show about sensitivity to a composed channel? The last question concerns the detector's response to specified operators. Evidence for an independent forensic cue would require another measurement.

\begin{samepage}
\subsection{Main Contributions}
The main contributions of this study are as follows.
\begin{itemize}
\item \textbf{Controlled evidence on training-view exposure.}
Two paired MobileNetV3Small--LSTM seed blocks quantify the effect of substituting sampled channel views for clean inputs under a common training and selection rule. A locked 1,088-recording evaluation tests the fitted detectors on a speaker-disjoint subset of the same benchmark release (Sections~\ref{sec:methods} and~\ref{sec:setup}).
\item \textbf{Class- and attack-specific robustness characterization.}
Stressed ROC-AUC improves in both blocks on development and held-out data. Misses increase on A05 and A06 in development and on A17--A19 in the held-out subset. Under noise followed by Opus, the 9207 channel model flags fewer REAL recordings and misses more FAKE recordings than ERM. Paired decisions identify the sources behind these trade-offs, and speaker-cluster intervals describe uncertainty within each fitted comparison (Section~\ref{sec:results}).
\item \textbf{Source-linked transformation-order diagnostics.}
Noise-to-Opus and Opus-to-noise views of the same recording support directional score contrasts and classwise threshold transitions. Attack-level analysis and a single-seed matched-component nuisance diagnostic delimit what can be attributed to the tested processing chains (Section~\ref{sec:results}). The contrasts describe detector responses and leave the underlying forensic mechanism unresolved.
\end{itemize}
\end{samepage}
